% !PRINTED 606
\weberpage{606}
%<<CONT
and when one considers that the fundamental number $\Delta$ is equal to $4d$ in case a), and equal to $d$ in case b), then, if we set
\weq{5}{x = 2a_{1,2} \text{ or } {}= 2a_{1,2}+1}
we obtain the following result:

1. \gesperrt{The necessary and sufficient condition for $p$ to contain two prime factors is that the congruence}
\weq{6}{x^2 \equiv \Delta\ (\mathrm{mod}\ 4p)}
\gesperrt{can be satisfied by an integral rational number $x$, or that the fundamental number $\Delta$ is a quadratic residue of $4p$.}

If we denote by $\pi, \pi'$ the two conjugate forms
\weq{7}{\pi = pu + (a_{1,2}+\Theta)v,\quad \pi' = pu + (a_{1,2}+\Theta')v,}
then multiplication yields
\weq{8}{\pi\pi' = p\left(pu^2 + xuv + \frac{x^2-\Delta}{4p}v^2\right).}

Now in this,
\[
pu^2 + xuv + \frac{x^2-\Delta}{4p}v^2
\]
is a unit, because, if $p$ does not divide $\Delta$, $x$ is not divisible by $p$, and if $p$ does divide $\Delta$, then $x$ is indeed divisible by $p$, but $(x^2-\Delta:4p)$ is not. For if $p$ is odd, $p^2$ cannot divide $\Delta$; and if $p=2$ and $\Delta$ is divisible by 2, then $\Delta = 4d$, $d \equiv 2, 3\ (\mathrm{mod}\ 4)$. But if we had
\[
x^2 - \Delta \equiv 0\ (\mathrm{mod}\ 16),
\]
it would follow that
\[
d \equiv \left(\frac{x}{2}\right)^2 (\mathrm{mod}\ 4),
\]
whereas $\left(\dfrac{x}{2}\right)^2$, being a square, cannot be congruent to 2 or 3. Accordingly $p$ is decomposed into the two prime factors $\pi$ and $\pi'$.

The two factors $\pi$ and $\pi'$ of $p$ can also be associated with one another. This occurs only when $\pi'$, and consequently also $\pi - \pi'$, is divisible by $\pi$. But now by (7)
\[
\pi - \pi' = (\Theta - \Theta')v,
\]
and consequently
\[
(\Theta - \Theta')^2 = \Delta
\]
%CONT>>
